%قبل از شروع مسئله ترجیحاً چیزی ننویسید. اما اگر خواستید بنویسید، در قالب لیستی (فقرات و شمارش) بنویسید.
%برای اضافه کردن نمره، بعدی از شروع‌مسئله، [x] را اضافه کنید و به جای x  نمره مورد نظر را قرار دهید. برای زیرمسئله‌ها هم به همین صورت است.
%\شروع‌مسئله[10]
\شروع‌مسئله
% صورت اصلی مسئله را در این‌جا بنویسید.
شنگدباو مسئول نقاشی دیواری شده که در ابتدای کار به رنگ سفید است. 
این دیوار به شکل یک مستطیل
$1 \times n$
است.
به او $q$ مرحله کار سپرده شده که باید به ترتیب انجام بشوند. در $i$امین مرحله او باید بازه 
$l_i$
تا
$r_i$
را به رنگی مخصوص مثل
$c_i$
در بیاورد. به‌عبارتی او باید هر مرحله مستطیلی
$1 \times (r_i - l_i)$
را به رنگ $c_i$ در بیاورد.

با توجه به اینکه ممکن است بازه‌هایی از دیوار قبلا به رنگ‌های متفاوتی در آمده‌باشند، او به ازای هر بخشی که از پیش رنگ شده و می‌خواهد رنگش را به
$c_i$
تغییر دهد قلمش را به رنگی در می‌آورد که ترکیبش با رنگ قبلی رنگ 
$c_i$
را بدهد. فرض کنید شنگدباو تمامی رنگ‌های مختلف را می‌تواند بسازد. به عبارتی اگر دیوارش به رنگ $x$ باشد و قرار باشد آن‌را به رنگ $y$ در بیاورد، او می‌داند باید چه رنگ $z$ ای را با $x$ ترکیب کند تا رنگ $z$ را بدهد و همیشه می‌تواند بخشی از دیوار را از رنگ $x$ به $z$ تغییر دهد. این تغییر رنگ قلم برای شنگدباو یک واحد خستگی به‌همراه می‌آورد.

در شکل
1
می‌توانید ۳ مرحله کار او را ببینید که نهایتا باعث می‌شود $5$ واحد خسته بشود.

\begin{figure}[H]
	\centering
	\includegraphics[scale=0.6]{figs/coloring}
	\label{fig:coloring}
	\caption{
در این شکل شنگدباو ابتدا بازه 
$2$ تا
$8$
را به رنگ قرمز در می‌آورد. سپس با استفاده از رنگی مرموز بازه $4$ تا $6$ را به رنگ زرد در می‌آورد. نهایتا باید بازه $3$ تا $7$ را مشکی کند. برای اینکار ابتدا بازه $3$ تا $4$ را با رنگ سبز ترکیب می‌کند. سپس بازه $4$ تا $6$ را با رنگ بنفش ترکیب می‌کند و سپس بازه $6$ تا $7$ را با رنگ سبز ترکیب می‌کند تا بازه 
$3$
تا
$7$
مشکی شود. هر سری که یک بازه‌را شنگدباو رنگ می‌کند یک واحد خسته می‌شود. لذا در مجموع ۵ واحد خسته می‌شود.	
}
\end{figure}

ابتدا نشان دهید میزان خستگی شنگدباو پس از $q$ مرحله از
$\mathcal{O}(q\times min(q, n))$
است. سپس با استفاده از تحلیل سرشکن شده نشان‌دهید درواقع میزان خستگی‌اش از
$\mathcal{O}(q)$
میشتر نخواهد شد.
%اگر مسئله چند قسمت داشت به این صورت عمل کنید:

\solution{
	پاسخ نمونه.
}




%%%%%%%%%%%%%%%%%%%%%%%%%%%%%%%%%%%%%%%%%%%%%%%%%%%%%
% برخی دستورات مفید برای تمرین‌ها. %
%%%%%%%%%%%%%%%%%%%%%%%%%%%%%%%%%%%%%%%%%%%%%%%%%%%%%

% فرمول‌های نمونه:
%$f(x) = \Theta(x)$
%$T(x, y) \leq \log(x) + 2^y + \binom{2n}{n}$
% {\tt code\_function}


% لیست موارد بدون شماره‌گذاری:
%\شروع{فقرات}
%    \فقره مورد اول 
%        \شروع{فقرات}
%            \فقره مورد اول داخلی 
%                \شروع{فقرات}
%                    \فقره اولین مورد خیلی داخلی 
%                    \فقره دومین مورد خیلی داخلی 
%                \پایان{فقرات}
%            \فقره مورد دوم داخلی 
%        \پایان{فقرات}
%    \فقره مورد دوم 
%\پایان{فقرات}

% لیست موارد شماره‌گذاری شده:
%\شروع{شمارش}
%    \فقره یک مورد 
%    \فقره یک مورد دیگر 
%    \فقره باز هم یک مورد دیگر
%\پایان{شمارش}

% برای نوشتن یک فرمول چند خطی که یک شماره بخوره. برای حالت‌های دیگر به سایت زیر مراجعه کنید:
% https://www.overleaf.com/learn/latex/Aligning_equations_with_amsmath
%\begin{align} 
%    \begin{split}
%        (x+y)^3 &= (x+y)^2(x+y) \\
%                &= (x^2+2xy+y^2)(x+y) \\
%                &= (x^3+2x^2y+xy^2) + (x^2y+2xy^2+y^3) \\
%                &= x^3+3x^2y+3xy^2+y^3
%    \end{split}                                 
%\end{align}

% نوشتن ماتریس و آرایه دوبعدی
%\begin{align}
%    A = 
%    \begin{bmatrix}
%        A_{11} & A_{21} \\
%        A_{21} & A_{22}
%    \end{bmatrix}
%\end{align}

% اضافه کردن تصویر با شرح. برای اضافه کردن تصویر بدون شرح، فقط دستور درج‌تصویر را استفاده کنید.
%\شروع{شکل}[h!]
%\تنظیم‌ازوسط
%  \درج‌تصویر[پهنا=0.5\پهنای‌متن]{img.jpg}
%\شرح{شرح تصویر}
%\برچسب{تصویر:۱}
%\پایان{شکل}
%\پایان{زیرمسائل}

%کد زیر را ببینید.
%%\lstinputlisting{source_filename.py}
%\begin{latin}
%\begin{lstlisting}
%# Import the modules
%import sys
%import random
%
%ans = True
%
%while ans:
%question = raw_input("Ask the magic 8 ball a question: (press enter to quit) ")
%
%answers = random.randint(1,8)
%
%if question == "":
%sys.exit()
%
%elif answers == 1:
%print "It is certain"
%
%elif answers == 2:
%print "Outlook good"
%
%elif answers == 3:
%print "You may rely on it"
%
%elif answers == 4:
%print "Ask again later"
%
%elif answers == 5:
%print "Concentrate and ask again"
%
%elif answers == 6:
%print "Reply hazy, try again"
%
%elif answers == 7:
%print "My reply is no"
%
%elif answers == 8:
%print "My sources say no"
%\end{lstlisting}
%\end{latin}
